% ==========================================================
% APRESENTAÇÃO - PRODUÇÃO AMBULATORIAL (SIA/SUS)
% IESB - Beamer Presentation (15 slides)
% ==========================================================

\documentclass[aspectratio=169,10pt]{beamer}

% ----------------------------------------------------------
% TEMA E CORES
% ----------------------------------------------------------
\usetheme{Madrid}
\usecolortheme{seahorse}

\definecolor{iesbBlue}{RGB}{27,79,114}
\definecolor{iesbLightBlue}{RGB}{46,134,193}
\definecolor{iesbRed}{RGB}{170,0,0}

\setbeamercolor{palette primary}{bg=iesbBlue,fg=white}
\setbeamercolor{palette secondary}{bg=iesbLightBlue,fg=white}
\setbeamercolor{palette tertiary}{bg=iesbBlue,fg=white}
\setbeamercolor{structure}{fg=iesbBlue}
\setbeamercolor{title}{fg=white,bg=iesbBlue}
\setbeamercolor{frametitle}{fg=white,bg=iesbBlue}
\setbeamercolor{block title}{bg=iesbBlue,fg=white}
\setbeamercolor{block body}{bg=iesbBlue!10}
\setbeamercolor{item}{fg=iesbBlue}

% Remover botões de navegação
\setbeamertemplate{navigation symbols}{}

% Número de slides
\setbeamertemplate{footline}{
    \leavevmode%
    \hbox{%
    \begin{beamercolorbox}[wd=.33\paperwidth,ht=2.25ex,dp=1ex,center]{author in head/foot}%
        \usebeamerfont{author in head/foot}\insertshortauthor
    \end{beamercolorbox}%
    \begin{beamercolorbox}[wd=.34\paperwidth,ht=2.25ex,dp=1ex,center]{title in head/foot}%
        \usebeamerfont{title in head/foot}\insertshorttitle
    \end{beamercolorbox}%
    \begin{beamercolorbox}[wd=.33\paperwidth,ht=2.25ex,dp=1ex,right]{date in head/foot}%
        \usebeamerfont{date in head/foot}\insertshortdate{}\hspace*{2em}
        \insertframenumber{} / \inserttotalframenumber\hspace*{2ex}
    \end{beamercolorbox}}%
    \vskip0pt%
}

% ----------------------------------------------------------
% PACOTES
% ----------------------------------------------------------
\usepackage[utf8]{inputenc}
\usepackage[T1]{fontenc}
\usepackage[brazil]{babel}
\usepackage{graphicx}
\usepackage{booktabs}
\usepackage{listings}
\usepackage{tikz}
\usepackage{fontawesome5}
\usepackage{amsmath}

\lstset{
    basicstyle=\ttfamily\tiny,
    frame=single,
    breaklines=true,
    backgroundcolor=\color{gray!10}
}

% ----------------------------------------------------------
% DADOS
% ----------------------------------------------------------
\title[Produção Ambulatorial SIA/SUS]{Análise da Produção Ambulatorial do SUS}
\subtitle{Extração, Armazenamento e Visualização de Dados do SIA/SUS}
\author[Matheus L. Ribeiro]{Matheus Lima Ribeiro}
\institute[IESB]{
    Centro Universitário IESB\\
    Ciência de Dados e Inteligência Artificial\\
    \medskip
    \textit{Integração de Técnicas e Projeto de Sistemas Inteligentes}\\
    Prof. Sérgio da Costa Côrtes
}
\date{Abril de 2026}

% ==========================================================
\begin{document}

% ----------------------------------------------------------
% SLIDE 1: TÍTULO
% ----------------------------------------------------------
\begin{frame}
\titlepage
\end{frame}

% ----------------------------------------------------------
% SLIDE 2: SUMÁRIO
% ----------------------------------------------------------
\begin{frame}{Sumário}
\tableofcontents
\end{frame}

% ----------------------------------------------------------
% SLIDE 3: INTRODUÇÃO - O SUS E O SIA
% ----------------------------------------------------------
\section{Introdução}
\begin{frame}{O SUS e o Sistema de Informações Ambulatoriais}

\begin{columns}[T]
\begin{column}{0.55\textwidth}
    \begin{itemize}
        \item \textbf{SUS}: Maior sistema público de saúde do mundo
        \item Atende $\sim$190 milhões de brasileiros
        \item \textbf{SIA/SUS}: Registra toda a produção ambulatorial
        \begin{itemize}
            \item Consultas, exames, terapias
            \item Procedimentos cirúrgicos ambulatoriais
        \end{itemize}
        \item \textbf{DATASUS/TabNet}: Portal de acesso aos dados
    \end{itemize}
\end{column}
\begin{column}{0.4\textwidth}
    \begin{block}{Dados do Projeto}
        \begin{itemize}
            \item \textbf{Período}: Jan/2024 -- Jan/2026
            \item \textbf{Meses}: 25
            \item \textbf{Abrangência}: Nacional
            \item \textbf{Granularidade}: Município
        \end{itemize}
    \end{block}
\end{column}
\end{columns}

\end{frame}

% ----------------------------------------------------------
% SLIDE 4: OBJETIVOS
% ----------------------------------------------------------
\begin{frame}{Objetivos do Projeto}

\begin{block}{Objetivo Geral}
Desenvolver um sistema integrado para extração, armazenamento e visualização de dados de Produção Ambulatorial do SIA/SUS.
\end{block}

\vspace{0.5cm}

\begin{columns}[T]
\begin{column}{0.48\textwidth}
    \begin{alertblock}{Objetivos Específicos}
    \begin{enumerate}
        \item Extrator HTTP puro (Requests + Regex)
        \item Banco de dados estruturado (SQLite)
        \item Aplicação web interativa (Streamlit)
    \end{enumerate}
    \end{alertblock}
\end{column}
\begin{column}{0.48\textwidth}
    \begin{exampleblock}{Entregas}
    \begin{enumerate}
        \item[4.] Estatísticas descritivas e gráficos
        \item[5.] Relatório técnico (LaTeX)
        \item[6.] Apresentação (Beamer)
    \end{enumerate}
    \end{exampleblock}
\end{column}
\end{columns}

\end{frame}

% ----------------------------------------------------------
% SLIDE 5: METODOLOGIA - ARQUITETURA
% ----------------------------------------------------------
\section{Metodologia}
\begin{frame}{Arquitetura do Sistema}

\begin{center}
\begin{tikzpicture}[
    box/.style={rectangle, draw=iesbBlue, fill=iesbBlue!10, thick, minimum height=1.2cm, minimum width=2.8cm, text centered, rounded corners=3pt, font=\small},
    arrow/.style={->, thick, iesbBlue}
]

% Nodes
\node[box, fill=iesbRed!15, draw=iesbRed] (tabnet) at (0,0) {DATASUS\\TabNet};
\node[box] (scraper) at (4,0) {Web Scraper\\(Python Requests)};
\node[box] (csv) at (8,0) {CSVs\\(Dados brutos)};
\node[box] (etl) at (8,-2.5) {Pipeline\\ETL};
\node[box] (db) at (4,-2.5) {SQLite\\(Banco)};
\node[box, fill=green!10, draw=green!50!black] (streamlit) at (0,-2.5) {Streamlit\\(Dashboard)};

% Arrows
\draw[arrow] (tabnet) -- (scraper);
\draw[arrow] (scraper) -- (csv);
\draw[arrow] (csv) -- (etl);
\draw[arrow] (etl) -- (db);
\draw[arrow] (db) -- (streamlit);

% Labels
\node[font=\tiny, above] at (2,0.1) {HTTP/POST};
\node[font=\tiny, above] at (6,0.1) {CSV (;)};
\node[font=\tiny, right] at (8.1,-1.25) {Pandas};
\node[font=\tiny, above] at (6,-2.4) {SQL};
\node[font=\tiny, above] at (2,-2.4) {Consultas};

\end{tikzpicture}
\end{center}

\vspace{0.3cm}
\begin{itemize}
    \item \textbf{3 camadas}: Extração $\rightarrow$ Armazenamento $\rightarrow$ Apresentação
    \item Pipeline automatizado e reprodutível
\end{itemize}

\end{frame}

% ----------------------------------------------------------
% SLIDE 6: METODOLOGIA - EXTRAÇÃO
% ----------------------------------------------------------
\begin{frame}{Metodologia: Robô de Extração}

\begin{columns}[T]
\begin{column}{0.55\textwidth}
    \textbf{Parâmetros da Extração:}
    \vspace{0.2cm}
    \begin{tabular}{ll}
    \toprule
    \textbf{Campo} & \textbf{Valor} \\
    \midrule
    Linha & Município \\
    Coluna & Subgrupo proced. \\
    Conteúdo & Qtd./Valor Aprovado \\
    Período & 25 meses \\
    Formato & CSV (;) \\
    \bottomrule
    \end{tabular}
\end{column}
\begin{column}{0.42\textwidth}
    \begin{block}{Engenharia e Resiliência}
    \begin{itemize}
        \item[\faRedo] Extração reversa de sessão do DATASUS
        \item[\faKey] Descoberta Regex dos tokens internos
        \item[\faShieldAlt] Re-encoding forçado ISO-8859-1 (Latin-1)
        \item[\faBolt] Zero timeouts vs Lentidão de browsers!
        \item[\faFileAlt] Headers Customizados do Client
    \end{itemize}
    \end{block}
\end{column}
\end{columns}

\vspace{0.3cm}
\begin{center}
\textbf{50 extrações} = 25 meses $\times$ 2 conteúdos (Quantidade + Valor)
\end{center}

\end{frame}

% ----------------------------------------------------------
% SLIDE 7: METODOLOGIA - BANCO DE DADOS
% ----------------------------------------------------------
\begin{frame}{Metodologia: Banco de Dados}

\begin{columns}[T]
\begin{column}{0.55\textwidth}
    \textbf{Tabela: producao\_ambulatorial}
    \vspace{0.2cm}

    \begin{tabular}{ll}
    \toprule
    \textbf{Campo} & \textbf{Tipo} \\
    \midrule
    id & INTEGER (PK) \\
    municipio & TEXT \\
    codigo\_municipio & TEXT \\
    uf & TEXT \\
    regiao & TEXT \\
    subgrupo\_proced. & TEXT \\
    periodo & TEXT \\
    qtd\_aprovada & INTEGER \\
    valor\_aprovado & REAL \\
    \bottomrule
    \end{tabular}
\end{column}
\begin{column}{0.42\textwidth}
    \begin{block}{Pipeline ETL}
    \begin{enumerate}
        \item Leitura dos CSVs
        \item Limpeza e normalização
        \item Enriquecimento (UF, Região)
        \item Carga no SQLite
        \item Consolidação Qtd+Valor
    \end{enumerate}
    \end{block}

    \vspace{0.3cm}
    \begin{alertblock}{Índices}
    município, uf, região, subgrupo, período
    \end{alertblock}
\end{column}
\end{columns}

\end{frame}

% ----------------------------------------------------------
% SLIDE 8: METODOLOGIA - STREAMLIT
% ----------------------------------------------------------
\begin{frame}{Metodologia: Aplicação Streamlit}

\begin{columns}[T]
\begin{column}{0.48\textwidth}
    \begin{block}{\faHome\ Dashboard}
    \begin{itemize}
        \item KPIs principais
        \item Evolução temporal
        \item Distribuição regional
        \item Filtros globais na sidebar
    \end{itemize}
    \end{block}

    \begin{block}{\faList\ Lista de Dados}
    \begin{itemize}
        \item Tabela interativa paginada
        \item Busca por município
        \item Exportação CSV
    \end{itemize}
    \end{block}
\end{column}
\begin{column}{0.48\textwidth}
    \begin{block}{\faChartBar\ Estatísticas}
    \begin{itemize}
        \item Tendência central
        \item Dispersão e quartis
        \item Histogramas
        \item Análise por dimensão
    \end{itemize}
    \end{block}

    \begin{block}{\faChartArea\ Gráficos (9 tipos)}
    Barras, Linhas, Pizza, Heatmap,
    Box Plot, Treemap, Dispersão,
    Área Empilhada, Barras Agrupadas
    \end{block}
\end{column}
\end{columns}

\end{frame}

% ----------------------------------------------------------
% SLIDES 9-12: RESULTADOS
% ----------------------------------------------------------
\section{Resultados}
\begin{frame}{Resultados: Visão Geral}

\begin{center}
\begin{tikzpicture}
    % KPI boxes
    \node[rectangle, draw=iesbBlue, fill=iesbBlue!10, rounded corners, minimum width=3.5cm, minimum height=1.8cm, text centered] (k1) at (0,0) {
        \textbf{\large 50}\\
        \small Extrações realizadas
    };
    \node[rectangle, draw=iesbLightBlue, fill=iesbLightBlue!10, rounded corners, minimum width=3.5cm, minimum height=1.8cm, text centered] (k2) at (4.5,0) {
        \textbf{\large 25}\\
        \small Meses de dados
    };
    \node[rectangle, draw=green!60!black, fill=green!10, rounded corners, minimum width=3.5cm, minimum height=1.8cm, text centered] (k3) at (9,0) {
        \textbf{\large 5.570+}\\
        \small Municípios
    };
\end{tikzpicture}
\end{center}

\vspace{0.5cm}

\begin{columns}[T]
\begin{column}{0.48\textwidth}
    \begin{block}{Extração}
    \begin{itemize}
        \item Quantidade Aprovada: \checkmark
        \item Valor Aprovado: \checkmark
        \item Período: Jan/2024 -- Jan/2026
        \item Formato: CSV (;)
    \end{itemize}
    \end{block}
\end{column}
\begin{column}{0.48\textwidth}
    \begin{block}{Banco de Dados}
    \begin{itemize}
        \item SQLite: Configuração zero
        \item Índices otimizados
        \item Views de agregação
        \item Metadados de extração
    \end{itemize}
    \end{block}
\end{column}
\end{columns}

\end{frame}

% ----------------------------------------------------------
\begin{frame}{Resultados: Distribuição Regional}

\begin{columns}[T]
\begin{column}{0.48\textwidth}
    \begin{block}{Concentração da Produção}
    \begin{itemize}
        \item \textbf{Sudeste}: Maior volume de produção ambulatorial, refletindo a concentração populacional e da rede assistencial
        \item \textbf{Nordeste}: Segunda maior produção
        \item \textbf{Sul}: Terceira posição
        \item \textbf{Centro-Oeste e Norte}: Menores volumes
    \end{itemize}
    \end{block}
\end{column}
\begin{column}{0.48\textwidth}
    \begin{alertblock}{Principais Observações}
    \begin{itemize}
        \item Distribuição correlacionada com população
        \item Disparidades regionais evidentes
        \item Diferenças nos subgrupos predominantes por região
        \item Variações temporais identificáveis
    \end{itemize}
    \end{alertblock}
\end{column}
\end{columns}

\end{frame}

% ----------------------------------------------------------
\begin{frame}{Resultados: Evolução Temporal}

\begin{block}{Análise Temporal (Jan/2024 -- Jan/2026)}
\begin{itemize}
    \item Tendência geral da quantidade aprovada ao longo do período
    \item Variações sazonais identificáveis (reduções em períodos de férias, feriados)
    \item Acompanhamento do valor aprovado permite análise financeira
    \item Dados recentes podem estar sujeitos a atualizações pelo Ministério da Saúde
\end{itemize}
\end{block}

\vspace{0.3cm}

\begin{exampleblock}{Tipos de Gráficos Temporais}
\begin{columns}[T]
\begin{column}{0.48\textwidth}
    \begin{itemize}
        \item Linha: evolução global
        \item Área empilhada: por região
    \end{itemize}
\end{column}
\begin{column}{0.48\textwidth}
    \begin{itemize}
        \item Barras: comparação mensal
        \item Tabelas: valores exatos
    \end{itemize}
\end{column}
\end{columns}
\end{exampleblock}

\end{frame}

% ----------------------------------------------------------
\begin{frame}{Resultados: Estatísticas Descritivas}

\begin{columns}[T]
\begin{column}{0.48\textwidth}
    \begin{block}{Medidas Calculadas}
    \begin{tabular}{ll}
    \toprule
    \textbf{Medida} & \textbf{Tipo} \\
    \midrule
    Média & Tendência central \\
    Mediana & Tendência central \\
    Moda & Tendência central \\
    Desvio Padrão & Dispersão \\
    Variância & Dispersão \\
    Amplitude & Dispersão \\
    Coef. Variação & Dispersão \\
    Q1, Q2, Q3 & Quartis \\
    IQR & Quartis \\
    Skewness & Forma \\
    Kurtosis & Forma \\
    \bottomrule
    \end{tabular}
    \end{block}
\end{column}
\begin{column}{0.48\textwidth}
    \begin{block}{Segmentações}
    Estatísticas calculadas para:
    \begin{itemize}
        \item Dados globais
        \item Por Região
        \item Por UF (Top 10)
        \item Por Subgrupo (Top 15)
        \item Por Período (temporal)
    \end{itemize}
    \end{block}

    \begin{alertblock}{Visualização}
    \begin{itemize}
        \item Histogramas com média/mediana
        \item Box plots por região
        \item Tabelas resumo
    \end{itemize}
    \end{alertblock}
\end{column}
\end{columns}

\end{frame}

% ----------------------------------------------------------
% SLIDE 13: BANCO DE DADOS
% ----------------------------------------------------------
\section{Banco de Dados}
\begin{frame}{Banco de Dados: Detalhes}

\begin{columns}[T]
\begin{column}{0.48\textwidth}
    \begin{block}{Estrutura}
    \begin{itemize}
        \item \textbf{Tabela principal}: producao\_ambulatorial
        \item \textbf{Metadados}: extracao\_metadata
        \item \textbf{Views}: vw\_producao\_por\_uf, vw\_producao\_por\_regiao
    \end{itemize}
    \end{block}

    \begin{exampleblock}{Consulta Exemplo}
    \begin{lstlisting}[language=SQL]
SELECT uf, regiao,
  SUM(quantidade_aprovada) as qtd,
  SUM(valor_aprovado) as valor
FROM producao_ambulatorial
GROUP BY uf, regiao
ORDER BY qtd DESC
LIMIT 10;
    \end{lstlisting}
    \end{exampleblock}
\end{column}
\begin{column}{0.48\textwidth}
    \begin{block}{Escolha do SQLite}
    \begin{itemize}
        \item[\faCheck] Configuração zero
        \item[\faCheck] Arquivo único portátil
        \item[\faCheck] SQL completo
        \item[\faCheck] Integração Python nativa
        \item[\faCheck] Performance adequada
    \end{itemize}
    \end{block}

    \begin{alertblock}{Índices Criados}
    \begin{itemize}
        \item municipio
        \item uf, regiao
        \item subgrupo\_procedimento
        \item periodo
        \item (uf, periodo) -- composto
    \end{itemize}
    \end{alertblock}
\end{column}
\end{columns}

\end{frame}

% ----------------------------------------------------------
% SLIDE 14: CONCLUSÕES
% ----------------------------------------------------------
\section{Conclusões}
\begin{frame}{Conclusões}

\begin{block}{Contribuições do Projeto}
\begin{enumerate}
    \item \textbf{Automação Reversa}: Extrator HTTP supera em 10x a lentidão de navegadores no TabNet.
    \item \textbf{Estruturação}: Pipeline ETL transforma dados brutos em base consultável
    \item \textbf{Acessibilidade}: Dashboard Streamlit democratiza o acesso à informação
    \item \textbf{Integração}: Projeto une web scraping, banco de dados, análise e visualização
\end{enumerate}
\end{block}

\vspace{0.3cm}

\begin{exampleblock}{Trabalhos Futuros}
\begin{columns}[T]
\begin{column}{0.48\textwidth}
\begin{itemize}
    \item Modelos de Machine Learning para previsão
    \item Migração para PostgreSQL
    \item Integração com SIH/SUS e CNES
\end{itemize}
\end{column}
\begin{column}{0.48\textwidth}
\begin{itemize}
    \item Deploy em nuvem
    \item Análise geoespacial
    \item Alertas automatizados
\end{itemize}
\end{column}
\end{columns}
\end{exampleblock}

\end{frame}

% ----------------------------------------------------------
% SLIDE 15: REFERÊNCIAS
% ----------------------------------------------------------
\section{Referências}
\begin{frame}{Referências}

\begin{thebibliography}{9}
\setbeamertemplate{bibliography item}[text]

\bibitem{datasus} DATASUS. \textit{Departamento de Informática do SUS}. Disponível em: \url{https://datasus.saude.gov.br}. Acesso em: abr. 2026.

\bibitem{siasus} MINISTÉRIO DA SAÚDE. \textit{Sistema de Informações Ambulatoriais do SUS (SIA/SUS)}. 2026.

\bibitem{tabnet} DATASUS. \textit{TabNet -- Informações de Saúde}. Disponível em: \url{https://datasus.saude.gov.br/informacoes-de-saude-tabnet/}. Acesso em: abr. 2026.

\bibitem{selenium} Selenium. \textit{Selenium Documentation}. Disponível em: \url{https://www.selenium.dev}. Acesso em: abr. 2026.

\bibitem{streamlit} Streamlit. \textit{Streamlit Documentation}. Disponível em: \url{https://docs.streamlit.io}. Acesso em: abr. 2026.

\bibitem{mitchell} MITCHELL, R. \textit{Web Scraping with Python}. O'Reilly Media, 2018.

\bibitem{paim} PAIM, J. et al. The Brazilian health system: history, advances, and challenges. \textit{The Lancet}, v. 377, p. 1778--1797, 2011.

\end{thebibliography}

\end{frame}

% ----------------------------------------------------------
\begin{frame}{}
\centering

\vspace{2cm}

{\Huge\color{iesbBlue}\textbf{Obrigado!}}

\vspace{1cm}

{\Large Matheus Lima Ribeiro}

\vspace{0.3cm}

{\normalsize Centro Universitário IESB\\
Ciência de Dados e Inteligência Artificial\\
Prof. Sérgio da Costa Côrtes}

\vspace{0.5cm}

{\small 2026/1}

\end{frame}

\end{document}
